\documentclass[11pt]{article}
\usepackage{mathblatt}
% gesetzt von werkzeuge/setzer.py; Lösungen nach bau/bauregeln.md „Lösungen“.
\newcommand{\kennung}{FP6}
\begin{document}
\pfheftstil
\fancyfoot[C]{{\footnotesize\color{mbgrau}\kennung\hfill\thepage}}
\pfheftkopf{Winkel an Geradenkreuzungen und Parallelen}{Klasse 7 \textperiodcentered{} Lösungen}

\begin{pfloesung}{}
\lz{1b)}{$\beta = 140^\circ$, $\gamma = 40^\circ$, $\delta = 140^\circ$}{$\gamma = 40^\circ$ (Scheitelwinkel); $\beta = 180^\circ - 40^\circ = 140^\circ$ (Nebenwinkel); $\delta = 140^\circ$ (Scheitelwinkel von $\beta$)}{}
\end{pfloesung}
\begin{pfloesung}{}
\lz{2.}{$\alpha = 52^\circ$, $\gamma = 52^\circ$, $\delta = 128^\circ$}{$\alpha = 180^\circ - 128^\circ = 52^\circ$ (Nebenwinkel von $\beta$); $\gamma = 52^\circ$ (Scheitelwinkel von $\alpha$); $\delta = 128^\circ$ (Scheitelwinkel von $\beta$)}{}
\end{pfloesung}
\begin{pfloesung}{}
\lz{3.}{$72{,}5^\circ$, $107{,}5^\circ$ und $107{,}5^\circ$}{Gegenüber: $72{,}5^\circ$ (Scheitelwinkel); daneben: $180^\circ - 72{,}5^\circ = 107{,}5^\circ$ (Nebenwinkel), zweimal}{}
\end{pfloesung}
\begin{pfloesung}{}
\lz{4.}{Nein: die Klingen sind nur $34^\circ$ geöffnet.}{Nein; Griffe und Klingen sind zwei Geraden, die sich am Niet kreuzen. Der Winkel zwischen den Klingen liegt dem Winkel zwischen den Griffen gegenüber: Scheitelwinkel, also $34^\circ$. $34^\circ < 40^\circ$.}{}
\end{pfloesung}
\begin{pfloesung}{}
\lz{5b)}{$\beta = 43^\circ$}{Der Winkel unten links zwischen $g$ und $h$ ist Scheitelwinkel zu $74^\circ$, also $74^\circ$; $k$ teilt ihn: $\beta = 74^\circ - 31^\circ = 43^\circ$}{}
\end{pfloesung}
\begin{pfloesung}{}
\lz{6.}{$\beta = 38^\circ$}{Der Winkel unten links zwischen $g$ und $h$ ist Nebenwinkel zu $114^\circ$: $180^\circ - 114^\circ = 66^\circ$; $k$ teilt ihn: $\beta = 66^\circ - 28^\circ = 38^\circ$}{}
\end{pfloesung}
\begin{pfloesung}{}
\lz{7.}{Mit der anderen Straße: $58^\circ$, also $6^\circ$ mehr.}{Der Radweg liegt im stumpfen Winkel: $180^\circ - 70^\circ = 110^\circ$ (Nebenwinkel). Zur anderen Straße: $110^\circ - 52^\circ = 58^\circ$. $58^\circ - 52^\circ = 6^\circ$: Mit der waagerechten Straße bildet er den größeren Winkel, um $6^\circ$.}{}
\end{pfloesung}
\begin{pfloesung}{}
\lz{8b)}{$\alpha = 56^\circ$, $\beta = 56^\circ$}{$\alpha = 56^\circ$ (Stufenwinkel: gleiche Stelle an $h$ wie $56^\circ$ an $g$); $\beta = 56^\circ$ (Wechselwinkel: über Kreuz zwischen den Parallelen)}{}
\end{pfloesung}
\begin{pfloesung}{}
\lz{9.}{$\alpha = 132^\circ$}{Stufenwinkel zu $48^\circ$ an $h$: $48^\circ$ (oben rechts); $\alpha$ liegt daneben: $\alpha = 180^\circ - 48^\circ = 132^\circ$}{}
\end{pfloesung}
\begin{pfloesung}{}
\lz{10.}{Sicher: spitzer Winkel an beiden Schienen $62^\circ$.}{Spitzer Winkel an Schiene 1: $180^\circ - 118^\circ = 62^\circ$ (Nebenwinkel); an Schiene 2 ebenso, weil die Schienen parallel sind (Stufenwinkel); $62^\circ > 60^\circ$.}{}
\end{pfloesung}
\begin{pfloesung}{}
\lz{11b)}{$\beta = 122^\circ$, $\gamma = 58^\circ$, $\delta = 122^\circ$}{$\beta = 180^\circ - 58^\circ = 122^\circ$; $\gamma = 58^\circ$; $\delta = 180^\circ - 58^\circ = 122^\circ$}{}
\end{pfloesung}
\begin{pfloesung}{}
\lz{12.}{$\gamma = 101^\circ$, $\delta = 114^\circ$}{$\delta = 180^\circ - 66^\circ = 114^\circ$ (an der Seite $AD$); $\gamma = 180^\circ - 79^\circ = 101^\circ$ (an der Seite $BC$)}{}
\end{pfloesung}
\begin{pfloesung}{}
\lz{13.}{$\beta = 96^\circ$, $\delta = 83^\circ$}{$\beta = 180^\circ - 84^\circ = 96^\circ$ (an der Seite $AB$); $\delta = 180^\circ - 97^\circ = 83^\circ$ (an der Seite $CD$)}{}
\end{pfloesung}
\begin{pfloesung}{}
\lz{14.}{Nein: links $34^\circ$ (zu steil), rechts $28^\circ$.}{Nein; Krone und Boden sind parallel, links unten $180^\circ - 146^\circ = 34^\circ$, rechts unten $180^\circ - 152^\circ = 28^\circ$; links $34^\circ > 30^\circ$: zu steil.}{}
\end{pfloesung}
\begin{pfloesung}{}
\lz{15.}{$112^\circ$}{$\beta = 180^\circ - 68^\circ = 112^\circ$ (Nachbarwinkel)}{}
\end{pfloesung}
\begin{pfloesung}{}
\lz{16.}{$\beta = 38^\circ$}{Scheitelwinkel zu $67^\circ$: $67^\circ$; $\beta = 67^\circ - 29^\circ = 38^\circ$}{}
\end{pfloesung}
\begin{pfloesung}{}
\lz{17.}{$\alpha = 39^\circ$}{$\alpha$ ist Scheitelwinkel zu $39^\circ$: $\alpha = 39^\circ$; die $61^\circ$ bei $B$ werden nicht gebraucht.}{}
\end{pfloesung}
\begin{pfloesung}{}
\lz{18.}{$\alpha = 123^\circ$}{Stufenwinkel an $h$: $57^\circ$ (oben rechts); $\alpha$ liegt darunter: $\alpha = 180^\circ - 57^\circ = 123^\circ$}{}
\end{pfloesung}
\begin{pfloesung}{}
\lz{19.}{$65^\circ$: $\alpha$, $\gamma$; \ $115^\circ$: $\beta$, $\delta$}{$65^\circ$: $\alpha$ (Stufenwinkel) und $\gamma$ (Scheitelwinkel zum Stufenwinkel); $115^\circ$: $\beta$ und $\delta$, je $180^\circ - 65^\circ = 115^\circ$ (Nebenwinkel).}{}
\end{pfloesung}
\end{document}
